% ---
% id: meat-core/lem:kraus-freedom
% kind: lem
% title: "Kraus freedom in the branch maps"
% summary: >-
%   Kraus factors with the same Gram matrices give branch operators related
%   by a partial isometry commuting with the key pinching, so the trace power,
%   deficit and pulled-back gradient agree. A subspace containing the support
%   is the support once the compression to it is positive definite.
% prominence: standard
% status: implemented
% uses: [def:branch, def:maps, def:joint-sectors, thm:branch-formula, lem:cone-concavity, thm:gradient, cor:full-gradient, thm:choi, def:support-frame, def:untrusted-choi, foundations/def:pinching, foundations/lem:matrix-function-perturbation]
% mentions: [meat-core/frank-wolfe/req:certified-evaluation]
% assumes: []
% refs: []
% provenance: newly-derived
% implements: certiqs_key.solvers.frank_wolfe.certified:_deficit
% lean:
% ---

The Kraus operators of a \tref{branch map} carry the square roots
$\sqrt{M_x}\otimes\sqrt{F_v}$ (Definition~\ref{def:branch}), and in the Choi
form the frame $W_{\tau}$ as well. A square root of an exact positive
semidefinite matrix is irrational in general, and is computed only
approximately. The lemma shows that nothing the objective or its gradient
reads depends on that choice: any factor with the same Gram matrix gives the
same branch values and the same pulled-back gradient, through a partial
isometry that the key pinching cannot see. The lemma is stated for one pair
of factor families at a time, so it holds for every element an enclosure
leaves possible, each with its own factors. A certified evaluation may
therefore build the branch maps from factors whose rank is fixed: exact
factors, or, for an element $LWL^{\dagger}$ known through an enclosed factor
$L$ and an exact positive diagonal core $W$, the factor $W^{1/2}L^{\dagger}$,
whose Gram matrix is that element for every $L$ the enclosure allows, though
$L$ itself is not known (Requirement~\ref{req:certified-evaluation}).
Part~(iv) then
compresses a branch operator, or a pinched block of one, to its support,
where it is positive definite and the divided-difference map of
\xref{foundations/lem:matrix-function-perturbation} applies: a subspace known
to contain the support, such as the range of the factors, is the support as
soon as the compression to it is shown positive definite.

\begin{lemma}\label{lem:kraus-freedom}
Let $\alpha\in(1,\infty)$, $p\coloneqq1/\alpha$ and $\hat c\in\Chkey$, and let
$\mathcal{H}_{\mathrm{in}}$ be a finite-dimensional Hilbert space. For every
outcome pair $(x,v)$ of a \tref{pass class} of $\hat c$, that is
$a(x,v)\in\Kann_{\mathrm{pass}}(\hat c)$, let $S_{xv}$ and $S'_{xv}$ be linear
maps from $\mathcal{H}_{\mathrm{in}}$ into finite-dimensional Hilbert spaces
$\mathcal{K}$ and $\mathcal{K}'$ with
\begin{equation}\label{eq:kraus-gram}
  S'^{\dagger}_{xv}S'_{xv}\;=\;S^{\dagger}_{xv}S_{xv} .
\end{equation}
For $\kappa\in\Kann_{\mathrm{pass}}(\hat c)$ put
\begin{equation}\label{eq:kraus-factored}
  K_{\kappa}\coloneqq\sum_{(x,v)\in a^{-1}(\kappa)}\ket{r(x,v)}_R\ket{x,v}_{XY}
  \ket{\kappa}_{C_AC_B}\otimes S_{xv},
\end{equation}
a map into $\mathcal{B}\coloneqq\Hs{R\,XY\,C_AC_B}\otimes\mathcal{K}$, and
$K'_{\kappa}$ likewise with $S'_{xv}$, into
$\mathcal{B}'\coloneqq\Hs{R\,XY\,C_AC_B}\otimes\mathcal{K}'$. For every
positive semidefinite operator $\omega$ on $\mathcal{H}_{\mathrm{in}}$ put
$X(\omega)\coloneqq\sum_{\kappa}K_{\kappa}\omega K_{\kappa}^{\dagger}$ and
$X'(\omega)\coloneqq\sum_{\kappa}K'_{\kappa}\omega K'^{\dagger}_{\kappa}$,
the sums over $\Kann_{\mathrm{pass}}(\hat c)$. Let $\cZ$ and $\cZ'$ be the
\tref[key pinching]{key pinchings} on $\mathcal{B}$ and $\mathcal{B}'$, by
the projectors $Z_j=\dyad{j}_R\otimes\id$, $j\in\Zset$, and $T$ and $T'$ the
functionals \eqref{eq:Tcone} with these pinchings. For $X\succeq0$ write
$\nabla T(X)$ for the operator of \eqref{eq:gradT-kernel} with the sum over
the strictly positive eigenvalues of $X$ only, $Y=\cZ(X^{p})$ and
$0^{\alpha-1}\coloneqq0$, which is the split form of
Theorem~\ref{thm:gradient}\,(i), and $\nabla T'$ likewise.
\begin{enumerate}[label=(\roman*),leftmargin=2.4em]
\item \emph{(A partial isometry.)} Let $\mathcal{W}\subseteq\mathcal{B}$ be
  the span of the subspaces
  $\ket{r(x,v)}_R\ket{x,v}_{XY}\ket{a(x,v)}_{C_AC_B}\otimes
  \operatorname{ran}S_{xv}$. There is a linear map
  $\mathbf U:\mathcal{B}\to\mathcal{B}'$ that is isometric on $\mathcal{W}$
  and vanishes on $\mathcal{W}^{\perp}$, with $\mathbf UK_{\kappa}=K'_{\kappa}$
  for every $\kappa$ and $(Z_j\otimes\idon{\mathcal{K}'})\,\mathbf U
  =\mathbf U\,(Z_j\otimes\idon{\mathcal{K}})$ for every $j$; and
  $\mathcal{W}$ is invariant under every $Z_j\otimes\idon{\mathcal{K}}$.
  Hence $X'(\omega)=\mathbf UX(\omega)\mathbf U^{\dagger}$, with
  $\operatorname{supp}X(\omega)\subseteq\mathcal{W}$, for every
  $\omega\succeq0$.
\item \emph{(Values.)} For every $\omega\succeq0$:
  $T'(X'(\omega))=T(X(\omega))$, $\tr X'(\omega)=\tr X(\omega)$,
  $\operatorname{supp}X'(\omega)=\mathbf U\operatorname{supp}X(\omega)$ and
  $\operatorname{rank}X'(\omega)=\operatorname{rank}X(\omega)$.
\item \emph{(Gradients.)} For every $\omega\succeq0$, with $X\coloneqq
  X(\omega)$ and $X'\coloneqq X'(\omega)$,
  \begin{equation}\label{eq:kraus-gradient}
    \nabla T'(X')=\mathbf U\,\nabla T(X)\,\mathbf U^{\dagger},
    \qquad
    \sum_{\kappa}K'^{\dagger}_{\kappa}\big(\nabla T'(X')-\id\big)K'_{\kappa}
    =\sum_{\kappa}K^{\dagger}_{\kappa}\big(\nabla T(X)-\id\big)K_{\kappa}.
  \end{equation}
\item \emph{(Compression to the support.)} Let $A\succeq0$ be an operator
  on $\mathcal{B}$.
  \begin{enumerate}[label=(\alph*),leftmargin=2.2em]
  \item \emph{(Containment.)} For every $\omega\succeq0$,
    $\operatorname{supp}X(\omega)\subseteq\sum_{\kappa}\operatorname{ran}K_{\kappa}$;
    and $\operatorname{supp}\big((Z_j\otimes\idon{\mathcal{K}})A
    (Z_j\otimes\idon{\mathcal{K}})\big)\subseteq(Z_j\otimes\idon{\mathcal{K}})
    \operatorname{supp}A$ for every $j$; and
    $\operatorname{supp}A^{q}=\operatorname{supp}A$ for every $q>0$.
  \item \emph{(A floor proves equality.)} If $Q:\mathbb{C}^{m}\to\mathcal{B}$
    is an isometry with $\operatorname{ran}Q\supseteq\operatorname{supp}A$
    and $Q^{\dagger}AQ$ is positive definite, then
    $\operatorname{ran}Q=\operatorname{supp}A$.
  \item \emph{(The compression.)} If $Q:\mathbb{C}^{m}\to\mathcal{B}$ is an
    isometry with $\operatorname{ran}Q=\operatorname{supp}A$ and $m\ge1$, then
    $N\coloneqq Q^{\dagger}AQ$ is positive definite, $A=QNQ^{\dagger}$,
    $A^{q}=QN^{q}Q^{\dagger}$ for every $q>0$, and for every operator $H$ on
    $\mathcal{B}$ the divided-difference sum of \eqref{eq:face-dpow} at $A$,
    over the strictly positive eigenvalues of $A$, equals
    $Q\,\mathrm{D}_N[Q^{\dagger}HQ]\,Q^{\dagger}$, with $\mathrm{D}_N$ the
    map of \xref{foundations/lem:matrix-function-perturbation}\,(iii) for
    $f(x)=x^{p}$.
  \end{enumerate}
\item \emph{(The branch maps.)} With $\mathcal{H}_{\mathrm{in}}=\Hs{A\Qsys_s}$
  and $S_{xv}=\sqrt{M_x}\otimes\sqrt{F^{(s)}_v}$, the $K_{\kappa}$ are the
  $K^{(s)}_{\kappa}$ of Definition~\ref{def:branch} and
  $X(\omega)=\cG_{\hat c,s}(\omega)$. In \tref{trusted-source mode}, with
  $\mathcal{H}_{\mathrm{in}}=\Hs{A_0\Qsys_s}$, a support frame and
  $S_{xv}=\sqrt{M_x}\,W_{\tau}\otimes\sqrt{F^{(s)}_v}$, the $K_{\kappa}$ are
  the $K^{(s)}_{\kappa}(W_{\tau}\otimes\idon{\Qsys_s})$,
  $X(J_s)=\cG^{J}_{\hat c,s}(J_s)$, and
  $\sum_{\kappa}K^{\dagger}_{\kappa}(\nabla T(X)-\id)K_{\kappa}
  =\chi^{\dagger}_{\Qsys_s}\big(\cG^{\dagger}_{\hat c,s}[\nabla T(X)-\id]\big)$,
  the operator of \eqref{eq:grad-pieces}; the same holds on a
  \tref{joint sector}
  $(t,s)$ with $M^{(t)}_x$ and the frame of the Alice sector $t$
  (Definition~\ref{def:joint-sectors}). In each case
  $S^{\dagger}_{xv}S_{xv}$ is $M_x\otimes F^{(s)}_v$, respectively
  $(W_{\tau}^{\dagger}M_xW_{\tau})\otimes F^{(s)}_v$. So for any $S'_{xv}$
  with that Gram matrix, the deficit $f_{\hat c,s}$, the branch operator's
  rank and the operator $\Delta^{(s)}_{\hat c}$ of \eqref{eq:grad-pieces}
  are given, by (ii) and (iii), by $X'$ in place of $X$.
\end{enumerate}
\end{lemma}

\begin{proof}
\emph{Isometric conjugation.} Let $\mathbf V$ be isometric on a subspace
$\mathcal{S}$ and vanish on $\mathcal{S}^{\perp}$, or be an isometry and
$\mathcal{S}$ its whole domain. Then $\mathbf V^{\dagger}\mathbf V$ is the
orthogonal projector onto $\mathcal{S}$: for $\phi$, $\psi$ with components
$\phi_{\mathcal{S}}$, $\psi_{\mathcal{S}}$ in $\mathcal{S}$,
$\langle\phi,\mathbf V^{\dagger}\mathbf V\psi\rangle
=\langle\mathbf V\phi_{\mathcal{S}},\mathbf V\psi_{\mathcal{S}}\rangle
=\langle\phi_{\mathcal{S}},\psi_{\mathcal{S}}\rangle$, an isometry preserving
inner products by polarisation. If $A\succeq0$ has
$\operatorname{supp}A\subseteq\mathcal{S}$,
write $A=\sum_ka_ku_ku_k^{\dagger}$ with $a_k>0$ and orthonormal
$u_k\in\mathcal{S}$; the $\mathbf Vu_k$ are orthonormal, so
$\mathbf VA\mathbf V^{\dagger}=\sum_ka_k(\mathbf Vu_k)(\mathbf Vu_k)^{\dagger}$
is a spectral decomposition with the same positive eigenvalues, its spectral
projector for an eigenvalue $\lambda>0$ is $\mathbf VP\mathbf V^{\dagger}$
with $P$ that of $A$, its support is $\mathbf V\operatorname{supp}A$, and
$(\mathbf VA\mathbf V^{\dagger})^{q}=\mathbf VA^{q}\mathbf V^{\dagger}$ for
every $q>0$, with $0^{q}\coloneqq0$ as in \xref{foundations/def:pinching};
moreover $\tr[\mathbf VA\mathbf V^{\dagger}]=\tr[A\,\mathbf V^{\dagger}
\mathbf V]=\tr A$, since $\mathbf V^{\dagger}\mathbf V$ is the identity
on $\mathcal{S}\supseteq\operatorname{supp}A$.

\emph{(i).} For each pair $(x,v)$ define $V_{xv}$ on
$\operatorname{ran}S_{xv}$ by $V_{xv}(S_{xv}\phi)\coloneqq S'_{xv}\phi$.
By \eqref{eq:kraus-gram},
$\|S_{xv}\phi\|^{2}=\langle\phi,S^{\dagger}_{xv}S_{xv}\phi\rangle
=\|S'_{xv}\phi\|^{2}$, so $S_{xv}\phi=S_{xv}\phi'$ forces
$S'_{xv}(\phi-\phi')=0$: $V_{xv}$ is well defined, linear and isometric.
Extend it by $0$ on $(\operatorname{ran}S_{xv})^{\perp}$. Write
$e_{xv}\coloneqq\ket{r(x,v)}_R\ket{x,v}_{XY}\ket{a(x,v)}_{C_AC_B}$; these
unit vectors are mutually orthogonal for distinct $(x,v)$, their $XY$
components being. Put $\mathbf U\coloneqq\sum_{(x,v)}e_{xv}e_{xv}^{\dagger}
\otimes V_{xv}$. On $\mathcal{W}=\bigoplus_{(x,v)}e_{xv}\otimes
\operatorname{ran}S_{xv}$ it acts as $V_{xv}$ on each orthogonal summand,
with orthogonal images, so it is isometric there, and it vanishes on
$\mathcal{W}^{\perp}$. Then $\mathbf UK_{\kappa}
=\sum_{(x,v)\in a^{-1}(\kappa)}e_{xv}\otimes V_{xv}S_{xv}=K'_{\kappa}$. Since
$Z_je_{xv}e_{xv}^{\dagger}=e_{xv}e_{xv}^{\dagger}Z_j$, both being
$e_{xv}e_{xv}^{\dagger}$ if $r(x,v)=j$ and $0$ otherwise, $\mathbf U$
commutes with the pinching projectors as stated, and each $Z_j$ maps every
summand of $\mathcal{W}$ to itself or to $0$. Finally
$X'(\omega)=\sum_{\kappa}\mathbf UK_{\kappa}\omega K_{\kappa}^{\dagger}
\mathbf U^{\dagger}=\mathbf UX(\omega)\mathbf U^{\dagger}$, and
$\operatorname{supp}X(\omega)\subseteq\sum_{\kappa}\operatorname{ran}K_{\kappa}
\subseteq\mathcal{W}$.

\emph{(ii).} Put $X\coloneqq X(\omega)$. By isometric conjugation,
$X'^{p}=\mathbf UX^{p}\mathbf U^{\dagger}$, $\tr X'=\tr X$, and the support
and rank statements hold. By (i),
$\cZ'(\mathbf UA\mathbf U^{\dagger})=\sum_j(Z_j\otimes\id)\mathbf UA
\mathbf U^{\dagger}(Z_j\otimes\id)=\mathbf U\cZ(A)\mathbf U^{\dagger}$ for
every $A$, the adjoint of the commutation relation giving
$\mathbf U^{\dagger}(Z_j\otimes\id)=(Z_j\otimes\id)\mathbf U^{\dagger}$.
So $Y'\coloneqq\cZ'(X'^{p})=\mathbf UY\mathbf U^{\dagger}$ with
$Y\coloneqq\cZ(X^{p})\succeq0$, and $\operatorname{supp}Y\subseteq\mathcal{W}$,
$\operatorname{supp}X^{p}$ lying in $\mathcal{W}$ and $\mathcal{W}$ being
invariant under every $Z_j$. Isometric conjugation again gives
$T'(X')=\tr[(\mathbf UY\mathbf U^{\dagger})^{\alpha}]
=\tr[\mathbf UY^{\alpha}\mathbf U^{\dagger}]=\tr[Y^{\alpha}]=T(X)$.

\emph{(iii).} With $Y$ as in (ii), $Y'^{\alpha-1}=\mathbf UY^{\alpha-1}
\mathbf U^{\dagger}$ and so $\cZ'(Y'^{\alpha-1})=\mathbf U\cZ(Y^{\alpha-1})
\mathbf U^{\dagger}$. Let $X=\sum_i\lambda_iP_i$ over the distinct strictly
positive eigenvalues of $X$; the strictly positive eigenvalues of $X'$ are the
same, with spectral projectors $\mathbf UP_i\mathbf U^{\dagger}$. As
$\operatorname{ran}P_i\subseteq\mathcal{W}$, on which
$\mathbf U^{\dagger}\mathbf U$ acts as the identity,
$\mathbf UP_i\mathbf U^{\dagger}\,\mathbf U\cZ(Y^{\alpha-1})\mathbf U^{\dagger}
\,\mathbf UP_j\mathbf U^{\dagger}=\mathbf UP_i\cZ(Y^{\alpha-1})P_j
\mathbf U^{\dagger}$, and summing with the coefficients
$\alpha\,(\lambda_i^{p}-\lambda_j^{p})/(\lambda_i-\lambda_j)$, respectively
$\alpha p\lambda_i^{p-1}$, gives the first identity of
\eqref{eq:kraus-gradient}. For the second,
$K'^{\dagger}_{\kappa}AK'_{\kappa}=K_{\kappa}^{\dagger}\mathbf U^{\dagger}A
\mathbf UK_{\kappa}$; with $A=\mathbf U\nabla T(X)\mathbf U^{\dagger}$ and
with $A=\id$, and $\mathbf U^{\dagger}\mathbf U$ acting as the identity on
$\mathcal{W}\supseteq\operatorname{ran}K_{\kappa}\cup\operatorname{supp}\nabla T(X)$
(the latter inside $\operatorname{supp}X$), both terms reduce to those with
$K_{\kappa}$ and $\nabla T(X)$.

\emph{(iv)(a).} The support of a positive semidefinite operator is its
range. $\operatorname{ran}X(\omega)\subseteq\sum_{\kappa}
\operatorname{ran}(K_{\kappa}\omega K_{\kappa}^{\dagger})\subseteq
\sum_{\kappa}\operatorname{ran}K_{\kappa}$. With
$Z\coloneqq Z_j\otimes\idon{\mathcal{K}}$, $\operatorname{ran}(ZAZ)\subseteq
Z\operatorname{ran}(A)$, since $ZAZw=Z(AZw)$. With
$A=\sum_ka_ku_ku_k^{\dagger}$, $a_k>0$ and orthonormal $u_k$, both $A$ and
$A^{q}=\sum_ka_k^{q}u_ku_k^{\dagger}$ have range the span of the $u_k$.

\emph{(iv)(b).} For $y\ne0$,
$(Qy)^{\dagger}A(Qy)=y^{\dagger}(Q^{\dagger}AQ)y>0$, so $Qy\notin\ker A$.
Hence $\operatorname{ran}Q\cap\ker A=\{0\}$ and
$m=\dim\operatorname{ran}Q\le\dim\mathcal{B}-\dim\ker A
=\operatorname{rank}A$. The inclusion
$\operatorname{supp}A\subseteq\operatorname{ran}Q$ gives
$\operatorname{rank}A\le m$, so the two subspaces have the same dimension and
are equal.

\emph{(iv)(c).} $QQ^{\dagger}$ is the orthogonal projector onto
$\operatorname{supp}A$, so $A=QQ^{\dagger}AQQ^{\dagger}=QNQ^{\dagger}$. For
$y\ne0$, $Qy\ne0$ lies in $\operatorname{supp}A$, where $A$ is positive
definite, so $y^{\dagger}Ny=(Qy)^{\dagger}A(Qy)>0$. Isometric conjugation
with $\mathbf V=Q$ and $N$ in the role of $A$ there gives
$A^{q}=QN^{q}Q^{\dagger}$ and shows that
the strictly positive eigenvalues of $A$ are the eigenvalues $\mu_a$ of $N$,
with spectral projectors $Q\Pi_aQ^{\dagger}$; so the sum of
\eqref{eq:face-dpow} is $\sum_{a,b}f^{[1]}(\mu_a,\mu_b)\,Q\Pi_aQ^{\dagger}H
Q\Pi_bQ^{\dagger}=Q\,\mathrm{D}_N[Q^{\dagger}HQ]\,Q^{\dagger}$ by the
definition of $\mathrm{D}_N$ in
\xref{foundations/lem:matrix-function-perturbation}\,(iii).

\emph{(v).} Substituting $S_{xv}$ in the definition of $K_{\kappa}$ gives
\eqref{eq:kraus} of Definition~\ref{def:branch}, and $X(\omega)$ is then
\eqref{eq:GTheta} for the sector. In the Choi form,
$(\sqrt{M_x}\otimes\sqrt{F^{(s)}_v})(W_{\tau}\otimes\id)
=\sqrt{M_x}W_{\tau}\otimes\sqrt{F^{(s)}_v}$, so the $K_{\kappa}$ are the
$K^{(s)}_{\kappa}(W_{\tau}\otimes\id)$, and $X(J_s)=\cG_{\hat c,s}
(\chi_{\Qsys_s}(J_s))=\cG^{J}_{\hat c,s}(J_s)$ by the definition of
$\chi_{\Qsys_s}$ in Theorem~\ref{thm:choi}. The Kraus adjoint
$\cG^{\dagger}_{\hat c,s}(A)=\sum_{\kappa}K^{(s)\dagger}_{\kappa}AK^{(s)}_{\kappa}$
and $\chi^{\dagger}_{\Qsys_s}(B)=(W_{\tau}^{\dagger}\otimes\id)B
(W_{\tau}\otimes\id)$ (Corollary~\ref{cor:full-gradient}) compose to
$\sum_{\kappa}K_{\kappa}^{\dagger}AK_{\kappa}$. On a joint sector the Kraus
operators act through $\sqrt{M^{(t)}_x}\otimes\sqrt{F^{(s)}_v}$
(\eqref{eq:joint-deficit}), and the same substitution applies. The Gram
matrices are $(\sqrt{M_x}\otimes\sqrt{F^{(s)}_v})^{2}=M_x\otimes F^{(s)}_v$
and $(\sqrt{M_x}W_{\tau}\otimes\sqrt{F^{(s)}_v})^{\dagger}
(\sqrt{M_x}W_{\tau}\otimes\sqrt{F^{(s)}_v})
=W_{\tau}^{\dagger}M_xW_{\tau}\otimes F^{(s)}_v$. The deficit
$f_{\hat c,s}=T\circ\cG_{\hat c,s}-\tr\circ\,\cG_{\hat c,s}$
(\eqref{eq:f-def}) is unchanged by (ii), and $\Delta^{(s)}_{\hat c}$, the
pulled-back operator of \eqref{eq:grad-pieces} with $\nabla T$ read in the
split form, is unchanged by (iii).
\end{proof}
